\section{Introduction}

The conversion between mechanical and other forms of energy is a ubiquitous process in many science and engineering applications.
Most notable from an energy footprint perspective is that a significant fraction of energy is inevitably lost by dissipating heat or vibrations (mechanical or acoustic).
Piezoelectric harvesting of mechanical into electric energy is therefore an interesting approach to reducing the energy
footprint of many applications that could for example help extend the battery life times of mobile devices or even
create self-powered electronic devices. 
Going in the opposite direction, generating mechanical deformations in response to electric stimuli
has multiple interesting applications such as for nano-actuators (e.g. piezo-motors in atomic force microscopy) or haptic feedback for microsurgery and other health related fields \cite{Dagdeviren2014}.
There is an ongoing interest to create organic piezoelectric materials to replace inorganic ceramics\cite{Werling0,Werling1, quan,Moody}.  
Organic materials offer several potential advantages, for example widely tunable mechanical properties [Matyjaszewski] and avoiding rare or
problematic elements such as lead or niobium.
Recently, polymer foam-based organic piezo-materials were reported with large piezoelectric responses (244 pC/N) \cite{Moody}, 
which demonstrates that it is possible to design organics that could replace inorganic ceramics. 

To realize the full potential of organic piezo-materials, it is necessary to develop
rational design approaches for organic piezoelectrics. 
Rationalizing organic polymer based piezoelectrics is very different from crystalline (bulk) materials in that it requires 
a multi-scale approach ranging from the (single) molecule limit to the bulk material. 
%A theory or even language to connect the molecular-level piezoelectric response
%with the bulk observables has not been presented before.  
However, there is no established formalism, let alone an established language, for {\em molecular} piezoelectricity. 
For example, the existing theory of piezoelectricity is based on notions from continuum mechanics (bulk scale), to define quantities such as strain and stress,
typically derived for crystalline (periodic) systems based on unit cells and their deformations. 
This approach is not able to describe the piezoelectric response of individual molecules and, consequently, their
connection to the bulk response. 
Developing a molecular understanding of piezoelectric responses is, however, important
to describe (noncrystalline) polymers, as well as responses of small (e.g. nano) systems that are inhomogeneous,
anisotropic, and aperiodic, so that the bulk description is not yet applicable. 
Examples of the latter include the electromechanical response of biomolecules as well as nano-scale
machines.  
%In this publication, we therefore aim to develop a theory of
%molecular piezoelectric response to develop a bottom-up understanding of organic piezoelectric materials and to aid in
%the development of rational design approaches. 
In this publication, we aim to establish a formalism and a language for {\em molecular} piezoelectric responses.
Establishing such a theory of molecular piezoelectric response is essential to being able to develop a bottom-up
(molecular-scale) understanding of piezoelectric responses in molecular materials. 
Furthermore, we also underline differences between {\em molecular} piezoelectric response and bulk response, 
emphasizing the need for a theory that can be applied and compare potential organic piezoelectric systems.


%Organic piezoelectric materials - appeals
%We want a molecular understanding of piezoelectricity / molecular piezoelectrics because: 
%a) biomolecules = single-molecule response
%b) nano-scale machines
%c) ties in with mechanochemistry = modifying chemical reactivity and physicochemical properties using mechanical
%deformation.

In our previous publications, we developed a computational approach to predict piezoelectric responses 
based on the definition of the converse piezoelectric effect, the deformation of a system in
response to an applied field. 
This approach is most natural in aperiodic calculations, where it is straightforward to apply an
electric field perturbation to the Hamiltonian. 
We established a simple computational procedure where the piezo-coefficient \dtt{} is
estimated by calculating the geometric response while applying a finite electric field. 
We then improved on this approach by developing an analytical expression that calculates \dtt{} from the zero-field 
geometric Hessian and dipole moment derivative. 
This approach has several advantages over our first procedure, for example by avoiding the
finite-field calculations (which can be problematic because of electronic instability of molecules in finite electric
fields) \cite{Yamabe1977, raey, raey1,raey2} and by cutting down on the computational cost via requiring only zero-field calculations (as opposed to
calculations at several finite field strengths). 
We demonstrated that both approaches yield comparable results for
molecular systems and used these approaches to explore piezoelectric responses in hydrogen-bonded systems. 
In fact, we showed that hydrogen bonding in 2-methyl-4-nitroaniline (MNA) gives rise to significant
piezoelectricity \cite{Werling1}. 
We then explored several examples of hydrogen donor-acceptor systems to help establish a general
rationale for the construction of systems with large \dtt{}. 
We found that our analytical expression is also useful in explaining piezoelectric responses by showing that a large \dtt{} requires both a large 
inverse Hessian (large compliance of the bond) and large dipole moment derivative. 
Point out papers by other authors that build on our approach / explore similar systems. 

The previously developed mathematical model exploited several circumstances to optimize the approach for the description
of piezoelectric responses in hydrogen-bonded systems. 
%There were some implicit assumptions both in the model and reflected in the mathematical derivation. 
Firstly, we tacitly assumed that the hydrogen bond length maps onto the
deformation that one would observe at the bulk scale. 
In other words, we only took deformation along the H-bond axis into account.  
Secondly, we assumed that relaxation of the monomers within the field were negligible and would not
drastically affect the hydrogen-bond length \cite{Werling1}.  
We also exclude any anisotropic effects of the applied field on the hydrogen-bond length by
applying the field in only one direction.  
These assumptions allowed us to drastically reduce the dimensionality of the problem and map it to a single
variable $z$, the hydrogen-bond distance, which gave rise to the majority of
the piezoelectric response. 
These assumptions were backed by calculations that
show that the intramolecular piezoelectric response is at least an order of
magnitude smaller than the intermolecular response \cite{Werling1}.  
The mapping of the hydrogen bond-length to the \dtt{} works thanks to crystal packing effects that
confine the majority of the response to one plane. 
In these respects the potential of this model to describe the piezoelectric
effect for general types of organic crystals is rather limited indeed. 
Firstly, it is in general not guaranteed that there is a single, dominating direction of
largest piezoelectric response and even if there is, it might not be
straightforward to determine this direction a priori. 
In fact, in a sense molecules give rise to the most anisotropic cases imaginable compared to the
rather symmetric periodic cells.  
Therefore, we aim to extend our previous
approach by taking the full response of the nuclear coordinates into account.
% This would take care of possible anisotropy in the response and also does not
% make any assumptions about coordinates of largest response. 
This would reveal the full anisotropy in the response of the system and also does not make any assumptions about coordinates of the largest response.
% The axis and magnitude of the largest piezoelectric response can of course be obtained a
% posteriori by diagonalizing the calculated piezoelectric tensor. 
We will show that the axis and magnitude of of the largest piezoelectric response for a given pair of coordinates in a system can be obtained 
by diagonalizing a matrix derived from the yet to be presented piezoelectric matrix.  
Furthermore, we derive a useful formula to calculate the piezoelectric matrix and discuss its connection to continuum strain mechanics. 
We will also underline the difficulty and impracticality of calculating a full 3rd rank piezoelectric tensor for molecules or small systems
while outlining connections of continuum strain theory to molecular deformations.  
The method we present for calculating the piezoelectric matrix can be used to understand the deformation properties between any two bodies in a system
and has great applicability to rational design of organic piezoelectric materials.
% Another, advantage is that this generalized approach can be automated more
% easily because it does not require any user input regarding direction of the
% response. 
Another advantage is that this generalized approach can be automated and only requires that the user specifies the orientation of the 
system on which they wish to perform the calculation so that they may interpret the results relative to the systems orientation.
We discuss this and many other facets which users might find useful to understand and predict the response of molecular deformation in the presence
of fields.

%Taking into account the full relaxation of the system
%will lead to more accurate piezoelectric coefficients at a lower cost, which we can then easily automate for screening.  We could
%also apply such methods to single molecules, since the system can no longer be broken into individual monomers which we can separate.
%Thus we will devise a method to accomplish these goals in the sections to come.


The present publication is the first in a series aiming at the development of a bottom-up understanding of
electromechanical responses in aperiodic (inhomogeneous and anisotropic) systems. 
In this publication we will develop a description of the miscroscopic (single-molecule) electromechanical response. How to  
(Not supposed to be rigorous calculation of \dd{} to be compared to a bulk
measurement, but rather means of comparison for responses between individual
molecules. This is an area that does not have a language or formalism yet, but
requires one to quantify and rationalize. This derivation interfaces with
multiple relevant areas. 
Of course it would be extremely informative if one
could develop an understanding of how molecular piezoelectric response (which
we are treating here) maps onto the response as measured at the bulk scale.
Making this connection will be subject of future work. However, we expect that
developing this quantitative language of molecular piezoelectric response will
be enormously important for analyzing nano-scale (single-molecule) machines. It also 
connects well with the emerging area of mechanochemistry and we believe that
our approach will help lend mechanochemistry a quantitative language. 

This paper is organized as follows. In section II, we provide a brief background of strain theory and the linear
piezoelectric equation in the established language of bulk deformations. Since we feel that many chemists are unfamiliar
with continuum mechanics, we provide a more complete introduction to strain theory in appendix A and derivation of the Green strain tensor 
in the connection to longitudinal strain. Section III presents
our derivation of molecular piezoelectric response, which is split into five parts. 
In section III.A we simplify the equations for the field derivative of longitudinal strain that we introduce in the background section.
In section III.B we indrocue the piezoelectric matrix and make connections from bulk strain theory to 
molecular systems.
In section III.C we describe how to construct the displacement field derivative for molecules and the 
nuance of how to project out unwanted motions.
In section III.D we show how to calculate the piezoelectric matrix for molecules and some ways it can be used
to screen for ``good'' organic piezoelectric candidates.
In section III.E we describe a rank-4 piezoelectric tensor which for a molecule takes the place of the 
rank-3 piezoelectric tensor field in a bulk system.
At points necessary in section III, we discuss how deformation response in molecules is uniquely distinguished 
from bulk response and the inherent anisotropy and inability to ascribe a full piezoelectric system to molecular systems.
Section IV outlines the computational procedure for calculating molecular piezoelectric responses. 
In section V we present numerical results for molecules of interest to illustrate the concepts derived in the previous
sections. 
We conclude with a discussion and outlook in section VI. 
